\documentclass[11pt]{article}
\usepackage[utf8]{inputenc}
\usepackage{amsmath}
\usepackage{amsfonts}
\usepackage{amssymb}
\usepackage{bm}
\usepackage{geometry}
\usepackage{xcolor}
\usepackage{booktabs}
\usepackage{hyperref}

\geometry{margin=1.0in}

\title{Inferring Spatial Connectivity Decay Exponent\\
from Graph Topology Alone}
\author{Jan Balewski, NERSC}
\date{\today}

\begin{document}

\maketitle

\begin{abstract}
A spatially embedded neuronal network is constructed by sampling
connections with probability proportional to
$D_{ij}^{-\delta_{\mathrm{ker}}}$, where $D_{ij}$ is the Euclidean
distance between neurons $i$ and $j$ and
$\delta_{\mathrm{ker}}\in[0,3]$ is a real-valued decay exponent.
The limiting cases interpolate between uniform random connectivity
($\delta_{\mathrm{ker}}=0$) and strongly local, lattice-like wiring
($\delta_{\mathrm{ker}}\to 3$).  We ask: given only the binary
adjacency matrix~$G$ --- without access to neuron positions or
distances --- can one estimate~$\delta_{\mathrm{ker}}$?
We present five purely graph-theoretic diagnostics, each producing a
single scalar that varies monotonically with $\delta_{\mathrm{ker}}$
and therefore serves as a proxy for the hidden spatial decay.
The common principle is that large $\delta_{\mathrm{ker}}$ yields a
lattice-like graph (high local clustering, assortative degrees, low
spectral dimension), while small $\delta_{\mathrm{ker}}$ yields a
small-world or Erd\H{o}s--R\'enyi-like graph (diffuse neighborhoods,
weaker clustering, higher effective dimensionality).
\end{abstract}

\tableofcontents
\newpage

% ===============================================================
\section{Setup and Notation}
\label{sec:setup}
% ===============================================================

\subsection{The Generative Model}
\label{sec:gen_model}

Consider $N$ neurons placed on a regular two-dimensional grid within a
rectangular domain $[0,L]\times[0,H]$ with minimum spacing
$d_{\min}>0$.  After random subsampling of grid sites and
lexicographic sorting, each neuron~$j$ occupies a position
$\mathbf{P}_j=(x_j,y_j)$.  The pairwise Euclidean distance matrix is
$D_{ij}=\|\mathbf{P}_i-\mathbf{P}_j\|_2$.

Connections are drawn according to a distance-dependent affinity
kernel.  For every ordered pair $(i,j)$ with $i\neq j$, define the
affinity
\begin{equation}
  V_{ij} = D_{ij}^{-\delta_{\mathrm{ker}}},
  \qquad i\neq j;\qquad V_{ii}=0,
  \label{eq:affinity}
\end{equation}
where $\delta_{\mathrm{ker}}\in[0,3]$ is a real-valued \emph{decay
exponent} that controls how sharply connection probability falls off
with distance.  Because grid nodes are separated by at least
$d_{\min}>0$, the affinity is everywhere finite.

The exponent $\delta_{\mathrm{ker}}$ interpolates continuously between
qualitatively distinct connectivity regimes:

\begin{center}
\renewcommand{\arraystretch}{1.3}
\begin{tabular}{cl}
\toprule
$\delta_{\mathrm{ker}}$ & Character of the resulting graph \\
\midrule
$0$   & Erd\H{o}s--R\'enyi: connection probability independent of
        distance \\
$\lesssim 1$ & Weak spatial bias; long-range edges abundant;
        small-world-like \\
$\approx 1$  & Moderate locality; intermediate clustering \\
$\approx 2$  & Strong locality; neighborhood structure clearly visible \\
$\gtrsim 2$  & Very tight spatial coupling; lattice-like communities \\
$3$   & Extreme locality; connections almost exclusively between
        nearest neighbors \\
\bottomrule
\end{tabular}
\end{center}

\noindent
Each presynaptic neuron~$j$ independently samples an integer
out-degree $k_j\sim\mathrm{Uniform}(k_{\min},k_{\max})$ and draws
$k_j$ distinct postsynaptic targets $\mathcal{T}_j\subset
\{1,\dots,N\}\setminus\{j\}$ without replacement from the normalized
distribution
\begin{equation}
  p_{ij}
  = \frac{V_{ij}}{\displaystyle\sum_{i'\neq j}V_{i'j}}
  = \frac{D_{ij}^{-\delta_{\mathrm{ker}}}}
         {\displaystyle\sum_{i'\neq j}D_{i'j}^{-\delta_{\mathrm{ker}}}},
  \qquad i\neq j.
  \label{eq:conn_prob}
\end{equation}

\subsection{The Inference Problem}
\label{sec:inference_problem}

We are given only the weighted connectivity matrix
$A_{\mathrm{true}}\in\mathbb{R}^{N\times N}$. Since $A$-matrix represents interaction between  neurons  which can be either excitatory or inhibitory, the values of $A_{ij}$ can be either positive or negative.

Neuron positions
$\mathbf{P}$, pairwise distances~$D$, and the value of
$\delta_{\mathrm{ker}}$ are all \emph{unknown}.  From
$A_{\mathrm{true}}$ we extract the binary adjacency matrix
\begin{equation}
  G_{ij} =
  \begin{cases}
    1 & A_{\mathrm{true},ij}^{\mathrm{off}}\neq 0,\;i\neq j,\\
    0 & \text{otherwise},
  \end{cases}
  \label{eq:adjacency}
\end{equation}
where $A^{\mathrm{off}}$ denotes the off-diagonal part of
$A_{\mathrm{true}}$ (self-loops are excluded).  The goal is to
estimate the decay exponent $\delta_{\mathrm{ker}}$ from~$G$ alone, or
at minimum to construct scalar diagnostics that vary monotonically
with $\delta_{\mathrm{ker}}$ and therefore distinguish, e.g.,
$\delta_{\mathrm{ker}}=1$ from $\delta_{\mathrm{ker}}=2$.

The key qualitative insight is that as $\delta_{\mathrm{ker}}$
increases, the graph becomes more spatially clustered: edges are
shorter, neighborhoods overlap more, the graph looks more like a
lattice.  Conversely, as $\delta_{\mathrm{ker}}$ decreases toward
zero, edges span the full domain, neighborhoods are diffuse, and the
graph approaches an Erd\H{o}s--R\'enyi random graph.  Each of the five
methods below detects a different combinatorial or spectral
consequence of this transition.

\subsection{Additional Notation}

Throughout, we use the following derived quantities from $G$:
\begin{equation}
  k_i^{\mathrm{in}}  = \sum_{j} G_{ij}, \qquad
  k_j^{\mathrm{out}} = \sum_{i} G_{ij}, \qquad
  |\mathcal{E}| = \sum_{i,j} G_{ij},
\end{equation}
denoting the in-degree and out-degree of each neuron and the total
number of directed edges.  The symmetrized adjacency is
$\tilde{G}=(G+G^{\!\top})/2$ with degree matrix
$\tilde{D}_{ii}=\sum_j\tilde{G}_{ij}$.  Neighborhoods are defined as
$\mathcal{N}_{\mathrm{in}}(i)=\{j:G_{ij}=1\}$ and
$\mathcal{N}_{\mathrm{out}}(j)=\{i:G_{ij}=1\}$.


% ===============================================================
\section{Method 1: Degree Assortativity}
\label{sec:assortativity}
% ===============================================================

\subsection{Definition}

The Newman assortativity coefficient measures the Pearson correlation
between the degrees at either end of an edge.  For the directed graph
$G$, define the in-degree of the target and the in-degree of the
source for each edge $(i,j)\in\mathcal{E}$ (meaning $G_{ij}=1$,
i.e., neuron~$j$ drives neuron~$i$):
\begin{equation}
  r = \frac{
    |\mathcal{E}|^{-1}\!\displaystyle\sum_{(i,j)\in\mathcal{E}}
    k_i^{\mathrm{in}}\,k_j^{\mathrm{in}}
    \;-\;
    \left[|\mathcal{E}|^{-1}\!\displaystyle\sum_{(i,j)\in\mathcal{E}}
    k_i^{\mathrm{in}}\right]
    \left[|\mathcal{E}|^{-1}\!\displaystyle\sum_{(i,j)\in\mathcal{E}}
    k_j^{\mathrm{in}}\right]
  }{
    \sigma_{\mathrm{in,target}}\;\sigma_{\mathrm{in,source}}
  },
  \label{eq:assortativity}
\end{equation}
where $\sigma_{\mathrm{in,target}}$ and $\sigma_{\mathrm{in,source}}$ are the
standard deviations of $\{k_i^{\mathrm{in}}\}_{(i,j)\in\mathcal{E}}$
and $\{k_j^{\mathrm{in}}\}_{(i,j)\in\mathcal{E}}$ over all edges.

\subsection{Dependence on $\delta_{\mathrm{ker}}$}

When $\delta_{\mathrm{ker}}$ is large, each neuron's targets are
concentrated in a tight spatial neighborhood.  A neuron with high
in-degree sits in a locally dense region and receives input
predominantly from nearby neurons, which themselves receive input from
the same dense region --- hence they too have high in-degree.
High-degree nodes connect to high-degree nodes, producing
\emph{assortative} mixing ($r>0$).

As $\delta_{\mathrm{ker}}$ decreases, edges reach farther, a neuron's
targets span diverse spatial neighborhoods with diverse degree
profiles, and the degree--degree correlation weakens.  In the limit
$\delta_{\mathrm{ker}}\to 0$ (Erd\H{o}s--R\'enyi), degrees are
independent of position and $r\to 0$.

The assortativity coefficient thus increases monotonically with
$\delta_{\mathrm{ker}}$:
\begin{equation}
  \delta_{\mathrm{ker}}^{(1)} < \delta_{\mathrm{ker}}^{(2)}
  \quad\Longrightarrow\quad
  r^{(1)} \lesssim r^{(2)}.
  \label{eq:r_monotone}
\end{equation}


% ===============================================================
\section{Method 2: Common-Neighbor Overlap (Jaccard Index)}
\label{sec:jaccard}
% ===============================================================

\subsection{Definition}

For each directed edge $j\to i$ (i.e.\ $G_{ij}=1$), the Jaccard
overlap between the in-neighborhood of the target and the
out-neighborhood of the source is
\begin{equation}
  J_{ij}
  = \frac{|\mathcal{N}_{\mathrm{in}}(i)
          \cap \mathcal{N}_{\mathrm{out}}(j)|}
         {|\mathcal{N}_{\mathrm{in}}(i)
          \cup \mathcal{N}_{\mathrm{out}}(j)|}.
  \label{eq:jaccard}
\end{equation}
This answers a purely combinatorial question: \emph{if $j$ drives $i$,
how many of $j$'s other targets also drive~$i$?}  The mean Jaccard
index over all edges is
\begin{equation}
  \bar{J}
  = \frac{1}{|\mathcal{E}|}
    \sum_{(i,j)\in\mathcal{E}} J_{ij}.
  \label{eq:jaccard_mean}
\end{equation}

\subsection{Dependence on $\delta_{\mathrm{ker}}$}

The Jaccard index is the graph-theoretic analog of spatial locality
without ever defining locality explicitly.  When
$\delta_{\mathrm{ker}}$ is large, each neuron's fan-out lands in a
compact spatial neighborhood; two neurons connected by an edge are
necessarily close, so their neighborhoods overlap substantially and
$J_{ij}$ is large.  When $\delta_{\mathrm{ker}}$ is small, fan-out is
diffuse, common neighbors between any linked pair become rare, and
$\bar{J}$ shrinks.  In the Erd\H{o}s--R\'enyi limit
($\delta_{\mathrm{ker}}=0$) with edge density $p$, the expected
Jaccard index scales as $\bar{J}\sim p/(2-p)$, independent of
spatial structure.

The mean Jaccard index therefore increases monotonically with
$\delta_{\mathrm{ker}}$:
\begin{equation}
  \delta_{\mathrm{ker}}^{(1)} < \delta_{\mathrm{ker}}^{(2)}
  \quad\Longrightarrow\quad
  \bar{J}^{(1)} \lesssim \bar{J}^{(2)}.
  \label{eq:J_monotone}
\end{equation}

% ===============================================================
\section{Method 3: Cycle Closure Decay Rate}
\label{sec:cycles}
\sectionmark{Method 3}
% ===============================================================

\subsection{Definition}

Powers of the adjacency matrix count directed walks: $(G^k)_{ij}$ is
the number of distinct walks of length~$k$ from~$j$ to~$i$.  In
particular, $\operatorname{tr}(G^{k+1})$ counts closed walks
(directed cycles) of length $k{+}1$, while
$\mathbf{1}^{\!\top}G^k\mathbf{1}=\sum_{i,j}(G^k)_{ij}$ counts all
walks of length~$k$.  Define the \emph{$k$-step transitivity} as the
fraction of length-$k$ walks whose endpoints are connected by a
direct edge (thus closing into a $(k{+}1)$-cycle):
\begin{equation}
  \mathcal{T}_k
  = \frac{\operatorname{tr}(G^{k+1})}
         {\mathbf{1}^{\!\top}G^k\,\mathbf{1}},
  \qquad k=1,2,\dots
  \label{eq:transitivity_k}
\end{equation}
At $k=1$ this reduces (up to normalization) to the directed
clustering coefficient.  The \emph{cycle closure decay rate} is the
log-slope of $\mathcal{T}_k$ over a short range:
\begin{equation}
  \xi
  = -\frac{d\log\mathcal{T}_k}{dk}
  \bigg|_{k=2,\dots,K},
  \label{eq:xi}
\end{equation}
estimated by linear regression of $\log\mathcal{T}_k$ on~$k$ for
$k=2,\dots,K$ (typically $K=5$ or $6$).

\subsection{Dependence on $\delta_{\mathrm{ker}}$}

For a graph embedded in two dimensions, the existence of short cycles
is a signature of spatial locality: a walk that stays in a compact
region is far more likely to return to its origin than one that
wanders across the domain.  When $\delta_{\mathrm{ker}}$ is large,
neighborhoods are tight and densely interconnected, so walks of
moderate length close into cycles readily --- $\mathcal{T}_k$ remains
high across $k$ and $\xi$ is small.

When $\delta_{\mathrm{ker}}$ is small, long-range edges carry walks
far from their starting point.  A two-step walk $j\to i\to i'$ may
traverse the full domain, and the probability that $i'$ connects back
to~$j$ is no higher than the background edge density.  Cycle closure
becomes increasingly unlikely with each additional step, so
$\mathcal{T}_k$ drops steeply and $\xi$ is large.

The cycle closure decay rate therefore increases monotonically as
$\delta_{\mathrm{ker}}$ decreases:
\begin{equation}
  \delta_{\mathrm{ker}}^{(1)} < \delta_{\mathrm{ker}}^{(2)}
  \quad\Longrightarrow\quad
  \xi^{(1)} \gtrsim \xi^{(2)}.
  \label{eq:xi_monotone}
\end{equation}

% ===============================================================
% ===============================================================
\section{Method 4: Persistent Homology}
\label{sec:persistent_homology}
\sectionmark{Method 4}
% ===============================================================

This metric measures the summed lifetime of all loop features in the graph while sweeping from strong to weak edge weights.

We probe the topological complexity of the graph by computing its
persistent homology.  The symmetrized adjacency matrix
$\tilde{G}_{ij}=\tfrac{1}{2}(G_{ij}+G_{ji})$ removes any
directionality inherited from the original interaction matrix.  From
$\tilde{G}$ we define edge weights $w_{ij}=|\tilde{G}_{ij}|$
(interaction strength, irrespective of sign) and construct a filtered
Vietoris--Rips complex $\mathcal{R}(\epsilon)$: a simplex
$[v_0,\dots,v_k]$ is included whenever every pairwise weight exceeds
the threshold,
\begin{equation}
  [v_0,\dots,v_k]\in\mathcal{R}(\epsilon)
  \quad\Longleftrightarrow\quad
  w_{v_i v_j}\geq\epsilon
  \;\;\forall\;0\leq i<j\leq k.
  \label{eq:rips}
\end{equation}
Sweeping $\epsilon$ from $\max_{ij}w_{ij}$ down to zero produces a
nested filtration; the persistence diagram records the birth--death
pairs $(b_i,d_i)$ for each homological feature that appears and
subsequently merges or closes.

Two scalar summaries are extracted.  The \emph{total persistence} in
dimension~$k$,
\begin{equation}
  \Pi_k = \sum_{i}\,(d_i - b_i),
  \label{eq:total_persistence}
\end{equation}
measures the cumulative lifetime of all $k$-dimensional features.
The \emph{persistent Betti curve} $\beta_k(\epsilon)$ counts the
number of features alive at filtration value~$\epsilon$; its integral
$B_k = \int \beta_k(\epsilon)\,d\epsilon$ gives a single-number
summary of topological richness.

For graphs embedded in two dimensions, only $k=1$ (loops) carries 
meaningful geometric information; higher-dimensional features such 
as enclosed cavities ($k=2$) are not expected in a planar embedding, 
and their computation requires enumerating $(k{+}1)$-cliques, which 
becomes prohibitively expensive for large graphs.  We therefore 
restrict attention to $\Pi_1$.

For lattice-like graphs ($\delta_{\mathrm{ker}}\geq 2$), the
dimension-$1$ barcode shows persistent loops reflecting the regular 
mesh, giving moderate~$\Pi_1$.  In the Erd\H{o}s--R\'enyi limit
($\delta_{\mathrm{ker}}=0$), random long-range edges create many 
short-lived loops, producing a cloud of short bars in~$\Pi_1$ with 
no dominant persistent features.

$\Pi_1$ thus offers a purely topological diagnostic of the 
lattice-to-random transition, complementing the combinatorial and 
spectral methods of Sections~\ref{sec:assortativity}--\ref{sec:jaccard}.


% ===============================================================
\section{Summary and Comparison}
\label{sec:summary}
% ===============================================================

The four diagnostics are collected in
Table~\ref{tab:summary}.  Each uses only the adjacency
matrix~$G$ (or equivalently the zero pattern of
$A_{\mathrm{true}}^{\mathrm{off}}$) unless otherwise noted.  None
requires neuron positions, pairwise distances, or spike data.

\begin{table}[h]
\centering
\renewcommand{\arraystretch}{1.25}
\begin{tabular}{clcccc}
\toprule
\S & Method & Scalar & $\delta_{\mathrm{ker}}\!\uparrow$
  & Weights & Computational cost \\
\midrule
\ref{sec:assortativity}
  & Degree assortativity & $r$ & $\uparrow$
  & no & $O(|\mathcal{E}|)$ \\
\ref{sec:jaccard}
  & Jaccard overlap & $\bar{J}$ & $\uparrow$
  & no & $O(|\mathcal{E}|\,k_{\max})$ \\
\ref{sec:cycles}
  & Cycle closure decay & $\xi$ & $\downarrow$
  & no & $O(K\,N^3)$ or eigendecomp \\
\ref{sec:persistent_homology}
  & Persistent homology & $\Pi_1$ & $\Pi_1\!\uparrow$ & optional & $O(n_{\mathrm{simp}}^3)$ reduction \\
\bottomrule
\end{tabular}
\caption{Summary of the four topological diagnostics for the decay
  exponent $\delta_{\mathrm{ker}}$.
  The arrow in the fifth column indicates the direction of change as
  $\delta_{\mathrm{ker}}$ increases (graph becomes more spatially
  local).  The ``Weights'' column indicates whether the method uses
  the magnitudes $|A_{\mathrm{true},ij}|$ beyond the binary
  connectivity pattern.
  For persistent homology, $n_{\mathrm{simp}}$ denotes the number of
  simplices in the filtered complex.}
\label{tab:summary}
\end{table}

Methods~1--3 operate exclusively on the binary adjacency matrix~$G$
and are entirely insensitive to the magnitudes of the interaction
weights $A_{\mathrm{true},ij}$.  Degree assortativity and Jaccard
overlap count edges and shared neighbors; cycle closure
decay works with powers of~$G$, both of
which depend only on which entries are nonzero.

Persistent homology (Method~4) is the only diagnostic that can
optionally exploit weight magnitudes.  As formulated in
Section~\ref{sec:persistent_homology}, the filtration is built from
$w_{ij}=|\tilde{G}_{ij}|$, where $\tilde{G}$ is the binary
symmetrized adjacency.  Because $\tilde{G}_{ij}\in\{0,\tfrac{1}{2},1\}$,
this filtration has only two nontrivial thresholds and is
effectively binary.  A richer filtration is obtained by replacing the
edge weights with the actual interaction strengths,
$w_{ij}=\tfrac{1}{2}(|A_{\mathrm{true},ij}|+|A_{\mathrm{true},ji}|)$,
so that strongly coupled pairs enter the complex before weakly coupled
ones.  This weighted variant introduces a continuous filtration
parameter and may sharpen the topological signatures, at the cost of
depending on weight magnitudes that the other three methods ignore.

The four methods probe complementary facets of the same underlying
phenomenon.  Degree assortativity (Method~1) detects correlations at
the scale of single edges.  Jaccard overlap (Method~2) examines
two-hop neighborhood structure.  Cycle closure decay (Method~3)
interpolates between local and global scales through the
parameter~$k$.  Persistent homology (Method~4) captures the
topological skeleton of the graph --- its 
loop structure --- as a function of interaction strength, providing a
diagnostic that is invariant under graph isomorphism and insensitive
to node labeling.

For practical estimation of $\delta_{\mathrm{ker}}\in[0,3]$, one may
generate a calibration family of synthetic graphs at a grid of
$\delta_{\mathrm{ker}}$ values (with the same $N$, $k_{\min}$,
$k_{\max}$, and E/I ratio as the target network), compute all four
scalars for each, and fit monotone calibration curves.  Given an
observed graph, the four diagnostics are evaluated and mapped through
the calibration curves to produce four independent estimates of
$\delta_{\mathrm{ker}}$, which can be combined (e.g., by averaging or
by inverse-variance weighting) into a single consensus estimate.

\end{document}
